\section{Design} \label{sec:design}

% The aim of design: write like the larousse. High level design for someone to
% take and improve upon. 
%
% Idea is to motivate that this is well designed and well thought through.
% Key design decisions:
        
% ------------------------
% Design introduction: what is the design philosophy? why?

\af is a complex piece of systems software and therefore required a
careful design phase. In a systems security project, added code and complexity
represent more chances for vulnerabilities to appear. Therefore, it was
important to design \af to be \textbf{as minimal as possible} while fulfilling
the promise of the proposed fine-grained filtering.

To do this, we defined a \textbf{threat model} (\ref{subsec:threat-model}) and
let primary requirements (\ref{subsec:requirements}) follow. These primary
requirements are presented alongside \textit{corollary requirements}
(\ref{subsubsec:corollary-reqs}) and design assumptions
(\ref{subsec:assumptions}).

\afg architecture is visualised in Figure \ref{fig:arch-overview}. Each
section's function is motivated in turn and key design decisions are justified.
The project's nature also requires some auxiliary tooling, which is discussed
after \afg design is presented.

\subsection{Threat Model}\label{subsec:threat-model}

% TODO: be more precise about <exploitation>

We started by creating a \textbf{threat model} which lays out the situations and
pre-conditions where \af is effective at protecting an application.

Firstly, we assume that 1) a \ac{ta} has gained \ac{rce} privileges on the host
machine \footnote{We refer to the machine running \af as the \textit{host
}}. That is, the \ac{ta} can execute any code they wish on the host. This is one of the most serious breaches that can happen and often
lead to denial of service, data loss, or data theft.%TODO: cite RCE bad

The attacker aims to use their \ac{rce} exploit to \textbf{escalate privilege}. This
means we assume that 2) the \ac{ta} doesn't have root privileges on the host,
but will act to get them.

% TODO: cite swc contains TA to compromised compt.
\af assumes that 3) some form of \textbf{software compartmentalisation} is in
place on the host machine. This compartmentalisation confines the \ac{ta} to the
compromised compartment, preventing the \ac{ta} 
from branching to an area of the address space that has a different allowed set
of system calls.

For completeness, we also explicitly assume that the attacker does not have a
kernel-level exploit. \todo{explicitly say that this is out of scope}.

Defining the high level requirements for \af is now possible, as we can reason
about exactly \textbf{what \af must protect against}.

% Use/reference BPF threat model
% https://www.linuxfoundation.org/hubfs/eBPF/ControlPlane%20%E2%80%94%20eBPF%20Security%20Threat%20Model.pdf#page=11.14

\subsection{Requirements} \label{subsec:requirements}

\afg key feature is that it should detect where an application makes a 
\textbf{disallowed syscall} and intervene accordingly. 

This intervention should be \textbf{user configurable}: warning the user, killing 
the malicious process, and killing all protected processes should all be
possible. This gives the user flexibility to choose how to trade-off
availability and data integrity/confidentiality on a per-application basis.

We should see a \textbf{significant reduction} in the set of syscalls a \ac{ta}
can access after comprimising an application, and we want to do this
\textbf{without detrimentally impacting performance}.\todo{see if these
reqiurements can be made more concrete without giving numbers}

To summarise 

% TODO: chatgpt these for clarity
\begin{enumerate}
    \item \af should map a syscall whitelist to each file-backed region of a
        process's \ac{vma}
    \item \af should detect \textbf{where} in the \ac{vma} a syscall was made
        from and \textbf{intervene} if the syscall isn't whitelisted
    \item \afg intervention policy should be \textbf{user configurable}: issuing
        a warning, killing the malicious process, and killing all processes
        being filtered should all be possible
    \item \af should provide a \textbf{significant level of privilege} reduction
        when compared to a regular seccomp filter at reasonable runtime slowdown
\end{enumerate}

\subsection{Corollary Requirements} \label{subsubsec:corollary-reqs}

The primary requirements listed above leave some questions unanswered.
% TODO: cite devs don't know app syscalls
Developers don't know which system calls their application makes - let alone
which shared libraries are responsible for each syscall. Therefore, we also
provide \textbf{automated whitelist generation tooling} with \af. 

Syscall filtering literature hasn't defined a standard way to measure the
reduction in privilege that a particular filtering strategy provides. Therefore,
we look to propose our own metric to quantify the level of privilege reduction
afforded by \af when compared to seccomp.

\subsection{Assumptions}\label{subsec:assumptions}

There are some rare cases which require significant engineering effort to cover.
In these cases, we make some \textbf{simplifying assumptions}. This was done
primarily to avoid adding unnecessary code and complexity to \afg design. 
 
One of these is handling the case where \acg{libc} address space changes during
execution. This is technically possible, but is rarely seen in production.
Care needs to be taken to ensure that this assumption doesn't introduce a
security risk during the implementation. System calls made in the case where 
\ac{libc} has been changed should be treated as disallowed syscalls.

% XXX: this requirement is satisfied by IMPLEMENTATION DETAILS
%  - libc moved => libc calls will be treated as non-libc
%  - => check to see if in libc whitelist; no libc whitelist guaranteed
%  - (because of generation) => disallowed

We also assume that \ac{libc} is mapped contiguously in memory. Again, this is
true (almost) without exception and reduces the codebase size and complexity.
It is also important to note that this assumption does not introduce a security
risk due to the way that syscall site identification is implemented and shown in
Section \ref{sec:implementation}. \todo{make this forward reference more precise}
%  TODO: Cite that this is true (I don't know if it's true; and maybe trivially
%  exploitable

With the relevant background fully explained and the requirements defined, we
can move on to presenting the design of \afss.~ 

\subsection{Architecture} \label{subsec:arch}

\begin{figure}[h]
    \centering
    \includegraphics[width=0.8 \linewidth]{./diagrams/TODO.pdf} 
    \caption{High level overview of \afg architecture}
    \label{fig:arch-overview}
\end{figure}

\af is comprised of a (kernel-space) backend and a userspace frontend.
\ac{bpf} maps and ringbuffers are used to allow the backend and frontend to
communicate. The frontend is responsible for spawning the filtered application,
handling configuration, and \textit{attaching} the \ac{bpf} program to it's
tracepoint (discussed further in Section \ref{sec:implementation}).

The backend of \af is made up of a single \ac{bpf} program and \ac{bpf}
maps/ringbuffers to communicate with userspace.
The program is attached to a \textbf{raw tracepoint}, and runs every time
a system call is made - even ones not called by by the app being filtered.

\subsubsection{Determining Which Syscalls to Filter}

This means \af needs a way to identify which system calls were made by the
protected app. \acp{pid} of filtered processes
need to be stored in \iac{bpf} map. Then, on every system call, we check if the
called process's \ac{pid} is in the map of filtered processes. If not, ignore
the system call. Otherwise, we continue with the filtering.

% TODO: cite nginx forks to deal with applications
We also want to apply the filter to any child processes the
application might create. The reason for this is two-fold: a common workflow for
web servers (such as nginx) is to fork on receiving a request; the second is
that if a \ac{ta} could just call fork() to remove the syscall filter, we
wouldn't have proposed a very secure solution.

To do this, we can also check if the parent process's \ac{pid} is in the follow
map. If it is, we add the calling \ac{pid} to the follow map, and continue with
the filtering. So, we only ignore syscall if \textbf{neither the current nor
parent \ac{pid} exists in the follow map}.

Having decided to filter a syscall, we need to identify where in the process's
address space the syscall was made. The value of the \ac{pc} seems like a good
place to start, but doesn't help here.
Almost all system calls are made via \ac{libc} in a dynamically linked 
application, and so the \ac{pc} almost always points to \ac{libc} instead of the
library that actually made the syscall.

\subsubsection{Finding the Syscall Site}\label{subsubsec:find_syscall}

Therefore, \af must find the \textbf{first non-\ac{libc} return pointer in the
userspace stack}, and treat this address as the true syscall site. In other
words, we trace through the sequence of function calls that led to the system
calls and find the first call that wasn't made by libc.

For example consider a ``hello, world'' program in C. \texttt{printf()} is a 
function implemented in \ac{libc} which invokes the \texttt{write()} syscall
behind the scenes. This means the userspace stack will look something like:

% TODO: turn this into a diagram

%                             [STACK GROWS =>] 
% ------------------------------------------------------------------
% main (0x$a) => printf (0x$b) => ... => write() (0x$pc)
% ^^^^^^^^^^^^^| LIBC     | ------------ XXXXXXXXXXXXXXX
%              | BOUNDARY |              PC value (incorrect site)
% true syscall site

\todo{Add stack diagram for hello world}

% TODO: move hello world to appendix
%
% \begin{lstlisting}
%     # include <stdio.h>
%
%     int main() {
%         printf("hello, world");
%
%         return 0;
%     } 
% \end{lstlisting}
%

The addresses shown in the diagram are \textbf{return pointers}: these are the
memory address of the instructions which added a new frame to the stack. In this
context, it's helpful to think of return pointers as the \textbf{address of a function}.

The idea is that here we want to classify this syscall as having been made from
the \texttt{main} binary, not from \ac{libc}. We can do this by
\textit{unwinding the stack} until we find a return pointer which isn't within
\acg{libc} address space. This return pointer is then used as the syscall
invocation site.

\todo{Add diagram showing correct address with stack unwinding, incorrect
without}

\subsubsection{Finding the Calling Shared Object File}\label{subsubsec:find_so}

Having found the memory address of the function which made the system call, we
can now figure out which shared library made the system call. Doing this is a
question of checking the process's \ac{vma} to find which memory region the
function belongs to. 

% TODO: remove reference to vma_struct; could access via /proc/pid/maps
Each shared library exists in the \ac{vma} as a \texttt{file backed memory region}.
In-kernel, the \ac{vma} is implemented as a \texttt{vma\_struct} struct with a
(indirect) reference to the name of it's backing file. Thus, finding which
shared library ``owns'' the syscall site is a question of finding the
corresponding \texttt{vma\_struct} struct and reading the filename.

\subsubsection{To Kill or not to Kill?}

Knowing which shared object file made the syscall allows us to find the correct
syscall whitelist and intervene if need be.

To do this, we keep a mapping from shared library filename to a set of allowed
system calls. We use the previously calculated filename as a key to this map to
retrieve some (pre-configured) whitelists. If the syscall number is in the
whitelist, no action is taken and the \ac{bpf} tracepoint returns with code 0.

If the syscall isn't on the whitelist, \af will intervene. The intervention
policy is user-configurable and has to be read from a config map. Userspace is
informed that a syscall has tripped the filter as the offending \ac{pid} is
written to a ringbuffer. If the policy is set to kill the process (or all 
processes in the \ac{pid} filtering map), \af sends a kill signal to the 
offending \ac{pid}. 

When userspace reads \iac{pid} from the ringbuffer, it acts according to how it was
configured. A warning will be logged in warn mode, but in ``kill all'' mode, a
kill signal will be sent to each \ac{pid} in the filtering map. This has to be
done asynchronously (by userspace) as \ac{bpf} only allows you to send signals
to the calling process.

The whole process is summarised in pseudocode in
Listing~\ref{lst:syscall-filter}.

% TODO: format nicely; make sure figure doesn't interrupt listing; declare as
% figure/table/something with caption; add to List of ...

\begin{listing}[ht]
  \caption{Pseudocode showing what happens when a syscall is made}\label{lst:syscall-filter}
  \begin{minted}[fontsize=\small, frame=single, linenos]{c}
  on every syscall {
      // only filter syscalls made from filter app or its children
      pid, ppid = get_pid(), get_ppid()

      if (pid not in filter_map and ppid not in filter_map) {
          return 0;
      }

      // protect forks of protected processes
      if (pid not in filter_map) {
          insert(filter_map, pid)
      }

      syscall_nr = get_syscall_num()

      // unwind the stack to find the first non-libc return pointer
      rp = find_syscall_site()

      // find which shared library the rp points to
      calling_so = lookup_vma_file(rp)
      whitelist = get_whitelist(calling_so)

      if syscall_nr not in whitelist {
          // intervene based on user configuration
          intervene()
      }

      // otherwise, success!
      return 0;
  }
  \end{minted}
\end{listing}

This design is minimal, modular, and easy to reason about. It can be achieved
with a single \ac{bpf} program running on the \texttt{raw\_tp/sys\_enter}
tracepoint with two core maps: one for storing the \acp{pid} of filtered apps
and the other for mapping shared library names to syscall whitelists.

What's been laid out here is the interesting part of the application
- the core filtering logic. The frontend is much larger (in terms of LoC) than
the \ac{bpf} program, and also serves important functions in making the system
work.

\subsubsection{Configuring Whitelists and \ac{libc}}

The frontend is responsible for parsing the \af whitelists and loading them into
the whitelist map. We looked to build on existing standards for defining
syscall whitelists but found none suitable for our use case. Seccomp filters,
for example, are defined as a \iac{bpf} program - this isn't user friendly and
doesn't allow for mapping shared libraries to syscall filters. Therefore, we defined our
own standard: a very simple \ac{toml} file. An example is provided in Appendix
\todo{appendix example of TOML whitelist}.

The frontend is also responsible for supporting performance optimisations
discussed in Section \ref{sec:implementation}. One of these optimisations is
to keep track of \acp{libc} address space in a map, as we assume \acp{libc} won't
change. The frontend is reponsible for finding the filtered app's \ac{libc}
address range and storing it in a map.

Ensuring that frontend functionality such as parsing whitelists is correct is
trivial with unit tests. Testing \ac{bpf} however is much more difficult.
Writing tests is technically possible but requires mocking a lot of key system
resources: crucially for us, the stack and the \ac{vma}. Having to mock core
functionality of a system often leads to brittle tests which aren't effective -
therefore, we chose a different method of validation and evaluation.

\subsection{Validation and Evaluation}

\todo{make this more brief and move some detail to evaluation}

Validation of \af happens in two phases: first, we make sure that the \ac{bpf}
program we've written is not interfering with normal syscall functionality. Then,
we test that our filtering system works - that is, \af intervenes if and only if
a syscall that is not on the whitelist is made.

To make sure that \af does not interfere with normal syscall functioning, we
propose to use the \acg{ltp} syscall test suite. We will run the suite on our
testbench with and without \af enabled and make sure that the results are the
same.

Due to the problems with mocking and automating testing (expanded on in Section
\ref{sec:evaluation}), we will favour a manual approach to ensuring that the 
syscall filtering works as expected. Writing a comprehensive test harness in requires 
significant engineering effort: likely more than developing the solution. 

Therefore, we opted for validation via manual
verification. During development, we will write a set of programs with a stack 
that is easy to inspect with tools like \texttt{gdb}. We will use a small set of
syscalls and dynamically generate whitelists (as discussed in
\ref{subsubsec:additional-tooling}). Using a restricted set of syscalls allows
us to manually add and remove them from the whitelists and observe program
behaviour.

Evaluation will also involve two phases. First, we will investigate the degree
of privilege reduction that a developer can expect from using \af on an
application. To do this, we define a \textit{privilege metric} and compare how
privileged seccomp-filtered applications are vs \afss-filtered applications.
\todo{forward-ref privilege metric}.

We then look to evaluate how costly \af is at runtime. To do this, we will take
a broad array of benchmarks informed by modern syscall filtering literature and
run them with \af enabled. We will also run the same set of syscalls with a
seccomp filter and with no filter, enabling comparison with a state of the art
filtering system and with a baseline.

We also intend to find use cases where \af is not suitable due to performance
overhead. Therefore, we will use our preliminary findings to stress \af and
speculate about application characteristics which most impact performance.

To do this evaluation, we need additional tooling. As discussed, we need two
way to generate per-library syscall whitelists. We also need some way to turn
these whitelists into a simple seccomp filter. The design of these tools are
briefly discussed in Section \ref{subsubsec:additional-tooling}.

\subsection{Additional Tooling}\label{subsubsec:additional-tooling}

The first additional tool that we implemented and \textbf{built in to \af} was
the whitelist generator. This whitelist generator works via dynamic analysis,
which is not a gold standard for whitelist generation. For this use case,
however, we think this acceptable - whitelist generation is orthogonal to our
work and is its own field of research. We have implemented a generator here
primarily to aid development and evaluation, but more work is needed to create a
production-ready generator based on static analysis.

We implement the whitelist generator using a lot of the same functionalities of
\af. The key difference is that after identifying the syscall site, we write a
syscall number to the syscall site's map (recall that at this stage \af would
check to see if the syscall number was in the site's map and intervene if need
be). When the program finishes executing (or on Ctrl-C), the frontend reads the
contents of the bpf whitelist map and marshals the results to a \ac{toml} file.
The pseudocode is provided in Listing \ref{lst:generator-pseudocode} and shows
that the generation code makes heavy use of \afg functionality.

\begin{listing}[ht]
  \caption{Pseudocode showing how dynamic analysis based whitelist generation is
  implemented}\label{lst:generator-pseudocode}
  \begin{minted}[fontsize=\small, frame=single, linenos]{c}
  on every syscall {
      // only filter syscalls made from filter app or its children
      pid, ppid = get_pid(), get_ppid()

      if (pid not in filter_map and ppid not in filter_map) {
          return 0;
      }

      // protect forks of protected processes
      if (pid not in filter_map) {
          insert(filter_map, pid)
      }

      syscall_nr = get_syscall_num()

      // unwind the stack to find the first non-libc return pointer
      rp = find_syscall_site()

      // find which shared library the rp points to
      calling_so = lookup_vma_file(rp)

      // at this point, the generation code becomes different to addrfilter

      // create a new whitelist for the calling_so if it doesn't already exist
      if calling_so not in whitelist_map {
        insert_to_map(whitelist_map, calling_so, new(whitelist))
      }

      // write the syscall number into the calling_so's whitelist, and update
      // the map
      whitelist = get_whitelist(calling_so)
      write_to_whitelist(whitelist, syscall_nr)
      insert_to_map(whitelist_map, calling_so, whitelist)
        
      // success!
      return 0;
  }
  \end{minted}
\end{listing}

The other tool that we need for the evaluation is one to launch an application
with a seccomp filter enabled. The tool takes as arguments an \af syscall whitelist
and a path to an executable. It then parses the whitelist and generates a single
seccomp filter by taking the union of each shared library's allowed set of
syscalls present in the \af whitelist.

This section looked to motivate the minimalist design philosophy, threat model,
and requirements for \af. We designed how \af should behave on each system call,
and discussed the need for kernel-space and user-space code. We also looked
forward to evaluating our solution and designed required additional tooling.

With these evaluation tools defined, the next step is to turn our design into a
working solution. In Section \ref{sec:implementation}, we show how the raw 
tracepoint attachment, BPF map structures, stack‑unwinding logic and userspace 
frontend were implemented to realise the filtering, whitelist generation and 
seccomp‑wrapper utilities described above.
